\documentclass[11pt,a4paper]{article}
\usepackage[utf8]{inputenc}
\usepackage[T1]{fontenc}
\usepackage[margin=1in]{geometry}
\usepackage{xcolor}
\usepackage{hyperref}
\definecolor{primarycolor}{RGB}{33,37,41}
\definecolor{secondarycolor}{RGB}{0,102,204}
\hypersetup{colorlinks=true,linkcolor=secondarycolor,urlcolor=secondarycolor}
\pagestyle{empty}
\setlength{\parskip}{0.9em}
\setlength{\parindent}{0pt}

\begin{document}

Sehr geehrte Damen und Herren,

als Pflegehelferin mit Erfahrung in der stationären Altenpflege und fünf Jahren Berufserfahrung als Krankenschwester in Polen bewerbe ich mich bei der Sozialstation Dreisamtal als Pflegehelferin für Ihren ambulanten Pflegedienst. Die Grundpflege, Mobilisation und Unterstützung bei der Nahrungsaufnahme gehören zu meinem Arbeitsalltag; außerdem bringe ich die Qualifikation Behandlungspflege LG1 mit.

Im Umgang mit älteren Menschen, auch mit Demenz, arbeite ich ruhig und geduldig. In meiner derzeitigen Einrichtung leite ich wöchentlich eine Singgruppe für Bewohnerinnen und Bewohner mit Demenz. Zudem habe ich gemeinsam mit der Wohnbereichsleitung Maßnahmen zur Sturzprophylaxe eingeführt; im folgenden Halbjahr sank die Zahl der Stürze im Wohnbereich von 14 auf 8. Ihr Leitsatz „Pflege mit Zeit und Respekt“ und die festen Tourengruppen sprechen mich an, weil mir eine respektvolle und verlässliche Pflege wichtig ist.

Mit Vivendi PD dokumentiere ich sicher; MediFox habe ich bislang nicht genutzt. Einen Führerschein Klasse B besitze ich derzeit nicht. Die theoretische Prüfung habe ich bestanden; die praktische Prüfung ist für Januar 2027 vorgesehen.

Gern erläutere ich Ihnen in einem persönlichen Gespräch, wie meine Erfahrung und meine Kenntnisse Ihr Pflegeteam unterstützen können. Ich freue mich auf Ihre Rückmeldung.

Mit freundlichen Grüßen

Katarzyna Wiśniewska

\end{document}